\documentclass{situationist-preprint}

\title{Successful Intentions Only}
\msid{AES-09-a}
\msfamily{AES-09}
\msdiscipline{Aesthetics}
\mstype{Principal manuscript}
\mskeywords{intention; interpretation; intentional fallacy; actual intentionalism; hypothetical intentionalism; work meaning; irony}

\begin{document}
\maketitle

\begin{abstract}
Does the author's intention fix a work's meaning? Anti-intentionalists answer that the meaning of a work is fixed by its words and their public conventions, and that the author's intention is either unavailable or irrelevant. Extreme intentionalists answer that a work means whatever its author meant by it. We argue for a position between them: intentions matter, but only successful ones, and failed intentions belong to biography. We define the success of a semantic intention as its realization in the work, in the sense that the work's features, read under the conventions and in the context available to its intended audience, support the intended meaning and make it recognizable. We show that failed intentions cannot fix a work's meaning, because a work cannot mean what its features and conventions cannot support, and that successful intentions do fix it wherever convention and context leave more than one reading open, because a work is the product of an intentional act and its meaning is the meaning of that act. Irony, allusion and resolved ambiguity are the central cases. We then show that hypothetical intentionalism misassigns meaning in cases where an audience's best hypothesis diverges from a successful actual intention.
\end{abstract}

\section{The problem}

Swift's \emph{A Modest Proposal} recommends that the children of the Irish poor be sold as food for the rich. No competent reader takes the work to recommend this, and the explanation everyone gives is that Swift did not intend it: he meant the proposal ironically, as an indictment of the treatment of Ireland. Yet the words of the pamphlet are the words of a recommendation. The case suggests that the author's intention fixes what a work means. Other cases suggest the opposite. A poet who intends a line to express serenity, but who writes a line that every competent reader hears as menacing, has not thereby written a serene line; and an author who claims, years after publication, that a novel meant something that nothing in it supports does not thereby change what the novel means. The question is whether, and how far, the author's intention fixes a work's meaning.

The modern debate begins with \citet{WimsattBeardsley1946}, who argued that the design or intention of the author is neither available nor desirable as a standard for judging the success of a literary work: if the author succeeded, the work shows what was intended, and if the author failed, the intention is not evidence about the work. \citet{Beardsley1958} developed the argument into a general anti-intentionalism about meaning, on which a text means what its words mean under public conventions. \citet{Barthes1977} radicalized the rejection, arguing that the unity of a text lies in its destination, the reader, and not in its origin. \citet{Hirsch1967} replied that without the author's intended meaning there is no determinate meaning to interpret, and that validity in interpretation requires a fixed object, which only the author's will can provide. \citet{KnappMichaels1982} went further: meaning and intended meaning are the same thing, so that the debate between intentionalists and their opponents is confused, and marks without an author, such as a poem traced by waves on a beach, are not language at all.

A second generation of views took a middle course. \citet{Levinson1992} defended \emph{hypothetical intentionalism}: a work means what an ideal reader, informed about the work and its public context, would best hypothesize that the author intended, whether or not the author actually intended it. \citet{Carroll1992,Carroll2000}, \citet{Iseminger1996}, \citet{Stecker2003} and \citet{Livingston2005} defended versions of \emph{modest actual intentionalism}: the author's actual intention fixes the meaning of a work when, and only when, it meshes with what the work's features and conventions allow. We defend a version of the latter, and we give it the precise form of the title: only successful intentions fix meaning, and failed ones are facts about the author rather than about the work.

The result bears on other questions. The moderate moralism of AES-08 holds that a moral flaw in a work is an aesthetic flaw when it prevents the work's prescribed response, and it relies on an account of which responses a work prescribes. On our account, the prescriptions that count are those of successful intentions.

\section{Framework}

We restrict attention to the meaning of works whose medium is governed by conventions of meaning, centrally literary works, though the framework extends to depiction and to music with programmatic content. We write $W$ for a work, $A$ for its author and $p$ for a candidate meaning.

\begin{definition}[Work meaning]\label{d:wm}
The \emph{meaning} of a work is the content that a correct interpretation of the work attributes to it: what the work says, represents, implies or expresses, including non-literal content such as irony, allusion and implicature.
\end{definition}

\begin{definition}[Semantic intention]\label{d:si}
A \emph{semantic intention} of $A$ with respect to $W$ is an intention, held by $A$ in making $W$, that $W$ (or a part of it) mean $p$. It is to be distinguished from $A$'s further intentions, for example that $W$ be admired or that it change its readers' views, and from \emph{categorial} intentions that $W$ be taken as belonging to a genre or form \citep{Levinson1992}.
\end{definition}

\begin{definition}[Permissible reading]\label{d:perm}
A meaning $p$ is a \emph{permissible reading} of $W$ if $W$'s features, read under the linguistic and artistic conventions and in the context available to the audience for which $W$ was made, can support $p$: the conventions assign $p$ to those features, or $p$ is derivable from them by the pragmatic principles that govern non-literal meaning in that context.
\end{definition}

\begin{definition}[Success]\label{d:succ}
A semantic intention of $A$ that $W$ mean $p$ is \emph{successful} if and only if (i) $p$ is a permissible reading of $W$, and (ii) $W$'s features, in that context, provide its intended audience with sufficient grounds to recognize that $p$ is intended. A semantic intention that is not successful is \emph{failed}.
\end{definition}

Definition~\ref{d:succ} adapts the condition that \citet{Livingston2005} calls the meshing of intention and text, and its second clause reflects the Gricean point that non-natural meaning is meaning intended to be recognized as intended \citep{Grice1957}. Clause (ii) is not satisfied merely because the audience could, by chance, arrive at $p$; the work must give them reason to. Irony in \emph{A Modest Proposal} satisfies both clauses: the proposal is permissible as an ironic reading under the conventions of eighteenth-century satire, and the work's escalating callousness, its arithmetical treatment of infants and its closing protestation of disinterest signal to a contemporary reader that the speaker's voice is not the author's.

\begin{principle}[Conventional constraint]\label{p:conv}
A work cannot mean what is not a permissible reading of it.
\end{principle}

\begin{principle}[Underdetermination]\label{p:under}
For many works, convention and context make more than one reading permissible, and some of these readings are incompatible with one another.
\end{principle}

\begin{principle}[Agency]\label{p:agency}
A work is the product of an intentional act of its author, and when a successful semantic intention selects one reading among those permissible, the work means what the act was successfully performed to mean.
\end{principle}

Principle~\ref{p:conv} is the lesson of Humpty Dumpty, who held that when he used a word it meant just what he chose it to mean (Lewis Carroll, \emph{Through the Looking-Glass}, ch.~6). Humpty Dumpty's \enquote{glory} does not mean \enquote{a nice knock-down argument}, however firmly he intends it to, because nothing in the conventions or the context available to Alice supports that reading. The principle does not require that conventions be fixed in advance of every utterance: \citet{Davidson1986} showed that interpreters routinely track meanings that depart from prior conventions, such as malapropisms, when context gives them the means to do so. Definition~\ref{d:perm} includes what context makes available for this reason; what it excludes is meaning for which the audience has no route at all.

Principle~\ref{p:under} records a fact of literary practice. Words have several senses, syntax admits several parsings, an utterance can be literal or ironic, and a phrase may or may not allude to an earlier text. Principle~\ref{p:agency} records the conception of a work as an utterance, which intentionalists of every kind share: works are made by someone, for an audience, and their meaning is the meaning of a communicative act \citep{Carroll1992,Stecker2003}.

\section{The argument}

\begin{lemma}\label{l:failed}
A failed semantic intention does not fix the meaning of a work.
\end{lemma}

\begin{proof}
Suppose $A$ intends $W$ to mean $p$ and the intention fails. Then by Definition~\ref{d:succ} either $p$ is not a permissible reading of $W$, or it is permissible but $W$ provides no grounds for recognizing it as intended. In the first case, by Principle~\ref{p:conv}, $W$ does not mean $p$, and the intention does not fix the meaning. In the second case, $p$ is one of the readings the work could bear, but the work gives its audience no reason to prefer it to its rivals. What then selects among the permissible readings is not $A$'s intention, which has left no trace in the work, but whatever else bears on the choice: the other features of the work and the conventions of the genre. If those features favor a reading other than $p$, $W$ means that reading; if they favor none, $W$ is ambiguous. In neither case is its meaning fixed by the failed intention.
\end{proof}

\begin{lemma}\label{l:success}
If $A$ has a successful semantic intention that $W$ mean $p$, then $W$ means $p$.
\end{lemma}

\begin{proof}
By Definition~\ref{d:succ}, $p$ is a permissible reading of $W$ and $W$ provides grounds for recognizing it as intended. If $p$ is the only permissible reading, $W$ means $p$ by Principle~\ref{p:conv} and convention alone, and the intention coincides with what the work means. If other readings are permissible (Principle~\ref{p:under}), the successful intention selects $p$ among them, and by Principle~\ref{p:agency}, $W$ means $p$.
\end{proof}

\begin{theorem}\label{t:main}
The author's intentions fix a work's meaning when and only when they are successful. Failed intentions do not bear on what a work means; they are facts about its author.
\end{theorem}

\begin{proof}
The \enquote{when} direction is Lemma~\ref{l:success}, and the \enquote{only when} direction is Lemma~\ref{l:failed}. A failed intention is a mental state of the author that the work does not realize, so it is a fact about the author and not about the work.
\end{proof}

Theorem~\ref{t:main} does not say that intentions matter only when they are redundant. The cases in which successful intentions do work are precisely those that Principle~\ref{p:under} describes, and they are not marginal. \emph{A Modest Proposal} is permissibly read as sincere: there are documents of the period that propose comparably callous measures without irony. Swift's successful intention to be ironic is what makes the ironic reading the work's meaning rather than one of two meanings it might have had. The same holds of allusion: a phrase in a poem that echoes a line of Milton alludes to Milton if the poet intended it to and the echo is recognizable, and is a coincidence if the poet did not, although the words are the same in both cases. Convention fixes that the phrase could allude; only the intention fixes that it does.

\begin{corollary}\label{c:evidence}
Evidence external to a work, such as an author's letters or statements, is relevant to interpretation as evidence about which permissible reading the author successfully intended, and is irrelevant as evidence of intended meanings that are not permissible readings or that the work gives no grounds to recognize.
\end{corollary}

\begin{proof}
By Theorem~\ref{t:main}, the intentions that bear on meaning are successful ones. External evidence that $A$ intended $p$ bears on meaning only by bearing on whether $A$ had a successful intention to mean $p$. If $p$ is permissible and the work offers grounds for recognizing it, the external evidence may confirm that those grounds point where they seem to. If $p$ is not permissible, or the work offers no grounds, no evidence about $A$ makes the intention successful, and the evidence is biographical.
\end{proof}

Corollary~\ref{c:evidence} answers the evidential half of the argument of \citet{WimsattBeardsley1946}. They held that if an author succeeded, the work itself shows what was intended, so that external evidence is superfluous, and if the author failed, external evidence is irrelevant. The second half is correct and is Lemma~\ref{l:failed}. The first half is not, because the work shows what was intended only to readers who are already in a position to recognize the grounds it offers; external evidence can put a later reader, who lacks the original audience's context, in that position.

\section{Against hypothetical intentionalism}

Hypothetical intentionalism agrees that a work's meaning is the meaning of an utterance, but it locates that meaning in the best hypothesis an ideal member of the intended audience would form about the author's intention, using publicly available information about the work and its context and excluding private evidence such as diaries \citep{Levinson1992}. The view is designed to preserve the autonomy of the work while explaining irony and allusion.

\begin{proposition}\label{p:hyp}
Hypothetical intentionalism and the success account diverge, and where they diverge the success account assigns the correct meaning.
\end{proposition}

\begin{proof}
The views diverge in two kinds of case. In the first, an author has a successful intention to mean $p$, but the best public hypothesis is that the author meant some other permissible reading $q$, because the public evidence is misleading, for instance because the author's views were widely misreported. The work gives grounds for recognizing $p$, by clause (ii) of Definition~\ref{d:succ}, but the ideal reader's hypothesis, weighing the misreported views, settles on $q$. Hypothetical intentionalism assigns $q$, and so ascribes to the work a meaning that its author did not intend and that is not required by convention. Nothing then explains why $q$ rather than $p$ should be the work's meaning except an error about the author, and interpretive practice treats such cases, once the error is corrected, as cases in which the work had been misread. In the second kind of case, the best hypothesis is $p$ but the author did not intend $p$ and intended nothing among the permissible readings. Hypothetical intentionalism assigns $p$; the success account says the work is ambiguous or means what its other features favor. The success account here is the more modest and does not ascribe to the work a selection that no one made. In both cases the meaning hypothetical intentionalism assigns depends on a hypothesis about an act, and when the hypothesis is false the act it describes did not occur.
\end{proof}

\citet{Levinson1992} accepts consequences of the first kind and defends them by appeal to the public character of literature: a work belongs to its readers once published, and they are entitled to the best reading the public evidence supports. The reply concedes our point in another form. What readers are entitled to believe on the public evidence is one question; what the work means is another, and it is because the two can come apart that we speak of readers being justified but wrong.

\section{Objections and replies}

\begin{objection}[The intentional fallacy]\label{o:wb}
Intentions are private mental states, generally unavailable to readers. A theory on which they fix meaning makes interpretation depend on what cannot be known \citep{WimsattBeardsley1946}.
\end{objection}

\begin{reply}
Successful intentions are not unavailable: by clause (ii) of Definition~\ref{d:succ}, a successful intention is one the work gives its audience grounds to recognize. The intentions that fix meaning on our account are therefore the ones whose content is, in principle, accessible through the work and its context. Intentions that are wholly private, in the sense of leaving no trace in the work, are failed intentions, and Lemma~\ref{l:failed} denies them any authority. The objection is sound against extreme intentionalism and not against the success account.
\end{reply}

\begin{objection}[Meaning is intention]\label{o:km}
Meaning just is intended meaning. A text without an intending author, such as marks left by waves, has no meaning, so there is no meaning for an intention to fail to fix. Failed intentions are impossible \citep{KnappMichaels1982}.
\end{objection}

\begin{reply}
The identification proves too much. If meaning just is intended meaning, Humpty Dumpty's \enquote{glory} means \enquote{a nice knock-down argument}, and there is no such thing as a writer failing to say what she meant. But writers revise drafts precisely because the drafts do not say what they meant, and the revision is an improvement in what the text means, not in what the writer intends. The wave case shows, at most, that meaning requires an utterance, which Principle~\ref{p:agency} grants. It does not show that every intention behind an utterance is realized in it.
\end{reply}

\begin{objection}[Works mean more than their authors intended]\label{o:more}
Great works have meanings their authors did not intend, which later readers discover. If only intentions fix meaning, these meanings are excluded.
\end{objection}

\begin{reply}
Theorem~\ref{t:main} says that successful intentions fix meaning where they operate, not that only intended meanings are meanings. A permissible reading that no intention selected and that no successful intention excludes remains part of what the work can be read to mean, and a work may support more readings than its author foresaw. What the success account excludes is a meaning incompatible with the author's successful intention: a reading of \emph{A Modest Proposal} as a sincere recommendation is not a discovery but an error. Unintended meanings are admitted on the same footing as intended ones, as permissible readings; they acquire the status of the work's meaning only when nothing selects against them.
\end{reply}

\begin{objection}[Idleness]\label{o:idle}
Success is defined by permissibility and recognizability, which are matters of convention and context. Intention therefore does no work: whatever a successful intention selects, the features of the work already select.
\end{objection}

\begin{reply}
Recognizability is relative to a recognized intention. The features of \emph{A Modest Proposal} give a reader grounds to recognize an ironic intention; they do not make the work ironic in the absence of one. Had the same text been written in earnest by a projector of the period, it would have been a sincere proposal whose tone was, to a sensitive reader, accidentally grotesque. The features are evidence of the intention, and the intention is what the evidence is evidence of. The conditions of success say when an intention is realized in a work; they do not replace the intention.
\end{reply}

\section{What the result does not establish}

The argument does not settle how much context the intended audience must have for a reading to be permissible, and so does not resolve every hard case; it says what follows once that is settled. It does not apply to aspects of a work's significance that are not part of its meaning, such as its relevance to later readers or its place in the history of a form; these may be fixed by facts with no connection to intention. It does not decide whether categorial intentions, which determine the genre a work belongs to, are subject to the same success condition, though we think they are. Finally, it says nothing about the value of a failed intention as evidence about the author, which may be great: the classification of such intentions as biography is a classification, not a dismissal. The companion manuscript AES-09-b examines what interpretation can still do with them.

\section{Conclusion}

The author's intention fixes a work's meaning when it is successful: when the meaning intended is one the work can bear, under the conventions and context of its audience, and when the work gives that audience grounds to recognize it as intended. Successful intentions do real work wherever convention leaves readings open, as it does in irony, allusion and ambiguity. Failed intentions do not fix meaning, since a work cannot mean what it cannot support, and they belong to the biography of the author. The position keeps what is right in the anti-intentionalist argument and in the intentionalist reply, and it assigns meaning more accurately than hypothetical intentionalism, which substitutes a hypothesis about an act for the act.

\begin{thebibliography}{99}\small
\bibitem[Barthes(1977)]{Barthes1977} Barthes, R. (1977). The death of the author. In \emph{Image, Music, Text}. Trans. S. Heath. Fontana.
\bibitem[Beardsley(1958)]{Beardsley1958} Beardsley, M.~C. (1958). \emph{Aesthetics: Problems in the Philosophy of Criticism}. Harcourt, Brace.
\bibitem[Carroll(1992)]{Carroll1992} Carroll, N. (1992). Art, intention, and conversation. In G. Iseminger (Ed.), \emph{Intention and Interpretation}. Temple University Press.
\bibitem[Carroll(2000)]{Carroll2000} Carroll, N. (2000). Interpretation and intention: The debate between hypothetical and actual intentionalism. \emph{Metaphilosophy}, 31(1--2), 75--95.
\bibitem[Davidson(1986)]{Davidson1986} Davidson, D. (1986). A nice derangement of epitaphs. In E. LePore (Ed.), \emph{Truth and Interpretation: Perspectives on the Philosophy of Donald Davidson}. Blackwell.
\bibitem[Grice(1957)]{Grice1957} Grice, H.~P. (1957). Meaning. \emph{The Philosophical Review}, 66(3), 377--388.
\bibitem[Hirsch(1967)]{Hirsch1967} Hirsch, E.~D. (1967). \emph{Validity in Interpretation}. Yale University Press.
\bibitem[Iseminger(1996)]{Iseminger1996} Iseminger, G. (1996). Actual intentionalism vs. hypothetical intentionalism. \emph{The Journal of Aesthetics and Art Criticism}, 54(4), 319--326.
\bibitem[Knapp and Michaels(1982)]{KnappMichaels1982} Knapp, S. and Michaels, W.~B. (1982). Against theory. \emph{Critical Inquiry}, 8(4), 723--742.
\bibitem[Levinson(1992)]{Levinson1992} Levinson, J. (1992). Intention and interpretation: A last look. In G. Iseminger (Ed.), \emph{Intention and Interpretation}. Temple University Press.
\bibitem[Livingston(2005)]{Livingston2005} Livingston, P. (2005). \emph{Art and Intention: A Philosophical Study}. Clarendon Press.
\bibitem[Stecker(2003)]{Stecker2003} Stecker, R. (2003). \emph{Interpretation and Construction: Art, Speech, and the Law}. Blackwell.
\bibitem[Wimsatt and Beardsley(1946)]{WimsattBeardsley1946} Wimsatt, W.~K. and Beardsley, M.~C. (1946). The intentional fallacy. \emph{The Sewanee Review}, 54(3), 468--488.
\end{thebibliography}

\end{document}
